\begin{textsample}{POS Dim 3 – df – Score 24.00 – df/df\_inf\_9.txt}  \label{ex:f3_pos_121}
6.2. \textbf{Brincar} e Interagir Autora : Clara Batista Lima Na Educação \textbf{Infantil}, as aprendizagens ocorrem em meio às relações sociais, tendo em vista que, a partir delas, a \textbf{criança} interage tanto com \textbf{crianças} da mesma faixa \textbf{etária} e de outras idades quanto com os \textbf{adultos}, o que contribuirá efetivamente para seu desenvolvimento.
Ressalta -se que as interações se estabelecem nas relações sociais, desde o \textbf{nascimento}, por meio de comunicação gestual, corporal e verbal.
Constituem -se como possibilidades de \textbf{ouvir} o outro, de \textbf{conversar} e trocar experiências e de aprender coletiva e colaborativamente.
A maneira como as relações sociais acontecem, no âmbito da \textbf{instituição} de educação para a primeira \textbf{infância}, influencia na qualidade do processo de aprendizagem e desenvolvimento.
Em vista disso, o coletivo, a troca de experiência, a relação com os objetos, pessoas e os elementos sociais e culturais contribuem para a constituição de vínculos com o outro e com o conhecimento, a \textbf{curiosidade}, o espírito investigativo, criativo e imaginativo.
Nas interações que se estabelecem em uma educação cuidadosa, a unidade afeto-intelecto precisa se consolidar, pois a atividade intelectual envolve a afetividade intrinsecamente como ações indissociáveis presentes nos relacionamentos humanos.
Portanto, em meio às práticas educativas, é essencial a possibilidade de expressão das emoções e dos sentimentos, pois as pessoas envolvidas nessa prática educativa afetam e são afetadas ( VIGOTSKI, 2009 ).
A compreensão da \textbf{criança} como ser que pensa e sente simultaneamente pode mensurar a relevância da afetividade como parte integrante do processo de aprendizagem e desenvolvimento, o que deve pautar a reflexão sobre as interações estabelecidas na \textbf{instituição} de educação para a primeira \textbf{infância}.
Assim, é importante conhecer as preferências das \textbf{crianças}, a forma delas participarem nas atividades, seus \textbf{parceiros} elegidos para os diferentes tipos de tarefas e suas narrativas.
Essas observações e \textbf{percepções} podem ajudar o profissional da educação a reorganizar as atividades de modo mais adequado à realização dos propósitos \textbf{infantis} e das aprendizagens coletivamente trabalhadas.
As interações criança/criança são essenciais e merecem conquistar tempos e espaços no planejamento e nas atividades.
Os pressupostos teóricos do Currículo em Movimento do Distrito Federal não entendem o desenvolvimento como uma conquista individual, mas coletiva e que ocorre a partir do caminho de desenvolvimento de cada \textbf{criança}, em meio às relações sociais e culturais.
Nas relações interpessoais, intra e intergeracionais, com os objetos da cultura e com os saberes, a \textbf{criança} aprende, desenvolve -se e humaniza -se.
Outro aspecto importante traz -nos Kishimoto ( 2010 ) ao afirmar a necessidade de integrar a educação ao \textbf{cuidado} e à \textbf{brincadeira}, apresentando como elementos exigidos a(s) : - Interação com o docente ; - Interação com os \textbf{pares} ; - Interação com os \textbf{brinquedos} e materiais ; - Interação entre \textbf{criança} e ambiente ; - Interações ( relações ) entre a \textbf{instituição} que oferta Educação \textbf{Infantil}, a família e/ou responsáveis e a \textbf{criança}.
Segundo Kishimoto ( 2010, p. 01 ), “ a opção pelo \textbf{brincar} desde o início da educação \textbf{infantil} é o que garante a cidadania da \textbf{criança} e ações pedagógicas de maior qualidade ”. \textbf{Brincando}, a \textbf{criança} lança \textbf{mão} de variadas formas de expressão : gesticula, fala, desenha, imita, \textbf{brinca} com \textbf{sons}, canta, entre outras possibilidades. \textbf{Brincar} é condição de aprendizagem, desenvolvimento e, por desdobramento, de internalização das práticas sociais e culturais.
Para as \textbf{crianças}, \textbf{brincar} é algo muito sério, sendo uma de suas atividades principais.
Enfatiza -se que essa atividade não é a que ocupa mais tempo da \textbf{criança}, mas aquela que contribui de modo mais decisivo no processo de desenvolvimento \textbf{infantil} ( ELKONIN, 2012 ).
Segundo Vigotski ( 2008 ), a \textbf{brincadeira} cria a chamada zona de desenvolvimento iminente³, impulsionando a \textbf{criança} para além do estágio de desenvolvimento que ela já atingiu.
Para o autor, o \textbf{brincar} libera a \textbf{criança} das limitações do mundo real, permitindo que ela crie situações imaginárias. \textbf{Brincar} é uma ação simbólica essencialmente social, que depende das expectativas e convenções presentes na cultura.
Quando duas \textbf{crianças} \textbf{brincam} de ser um \textbf{bebê} e uma mãe, por exemplo, fazem uso da imaginação, mas, ao mesmo tempo, não podem se comportar de qualquer forma ; devem obedecer às regras do comportamento esperado para um \textbf{bebê} e uma mãe, dentro de sua cultura.
Caso não o façam, correm o risco de não serem compreendidas pelos companheiros de \textbf{brincadeira}.
Contudo, de acordo com a Psicologia Histórico-Cultural, ninguém nasce sabendo \textbf{brincar}.
A \textbf{brincadeira} emerge da vida em sociedade entre os seres humanos.
Aprende -se pelas interações com outras \textbf{crianças} e com \textbf{adultos}, pelo contato com objetos e materiais, pela observação de outrem, pela reprodução e recriação de \textbf{brincadeiras}, pelas oportunidades ofertadas para isso.
Aprende -se nas \textbf{instituições} de Educação \textbf{Infantil}, em casa e na sociedade, nas interações que se estabelecem entre os familiares e amigos.
As possibilidades de exploração do \textbf{brinquedo}, por exemplo, dependem da ação dos \textbf{adultos} e do que a \textbf{criança} incorpora dessa relação.
Diante de tudo isso, sugerem -se algumas perguntas que podem orientar a vivência da \textbf{brincadeira} no cotidiano da \textbf{instituição} de Educação \textbf{Infantil} : - Por que as \textbf{crianças} \textbf{brincam}? - \textbf{Brincar} é realmente importante na \textbf{instituição} de Educação \textbf{Infantil}? - Qual a relação entre \textbf{brincadeira}, aprendizagens e desenvolvimento? - De que maneira organizar e incentivar \textbf{brincadeiras} que quebrem os estereótipos de gênero e etnia? - Como articular as \textbf{brincadeiras} e interações com as experiências da comunidade? - Como preservar a memória cultural popular e vinculá -la às novas tecnologias? - Como observar, acompanhar e participar das \textbf{brincadeiras} para estabelecer vínculos e contribuir para o desenvolvimento da \textbf{criança}? - Como contribuir com a imaginação \textbf{infantil} instigando a criatividade, investigação, \textbf{curiosidade}? - É possível e desejável inserir atividades \textbf{lúdicas}, jogos e cantigas tradicionais no repertório contemporâneo da \textbf{brincadeira} \textbf{infantil}? ³ Segundo Zoia Prestes ( 2012 ), a tradução correta para o termo utilizado por Vigotski é “ zona de desenvolvimento iminente ”, pois o desenvolvimento está na iminência de acontecer.
Vigotski estabeleceu um pensamento dialético e, nesse sentido, não tratava de etapas que se sucedem de forma linear.
Permitir que as \textbf{crianças} sejam protagonistas nas ações de \textbf{brincar} não significa deixá -las sem a supervisão e orientação de \textbf{adultos}.
A \textbf{criança}, em todos os espaços e tempos da \textbf{instituição} de Educação \textbf{Infantil}, é o centro do planejamento curricular : mesmo quando \textbf{brinca} sozinha, o professor precisa ter um olhar \textbf{atento} ao que está acontecendo, observando as ações, indagações e conquistas que as \textbf{crianças} estabelecem por meio das \textbf{brincadeiras}.
Tal como ressalta Kishimoto ( 2010 ), > São inúmeras as experiências expressivas, corporais e \textbf{sensoriais} das \textbf{crianças} pelo \textbf{brincar}.
Não se podem planejar práticas pedagógicas sem conhecer a \textbf{criança}.
Cada uma é diferente de outra e tem preferências conforme sua singularidade.
Em qualquer agrupamento \textbf{infantil}, há \textbf{crianças} que estão mais avançadas, outras, em ritmos diferentes.
Dispor de um tempo mais longo, em ambientes com variedade de \textbf{brinquedos}, atende os diferentes ritmos das \textbf{crianças} e respeita a diversidade de seus interesses ( KISHIMOTO, 2010, p. 4 ).
A \textbf{brincadeira} deve se fazer presente nos \textbf{gestos} e nas diferentes formas de apresentação oral, nos \textbf{brinquedos} e jogos e nos exemplos habituais dados pelos profissionais da educação.
Ela também precisa guiar outras atividades, como troca de fraldas, banho, alimentação e escovação dos dentes, independentemente da faixa \textbf{etária}.
A \textbf{brincadeira}, como prática educativa, possibilita que as interações entre as \textbf{crianças} e seus \textbf{pares} e entre elas e os \textbf{adultos} se constituam como um instrumento de promoção da imaginação, da experimentação e da descoberta.

% matched lemmas: adulto, atento, bebê, brincadeira, brincar, brinquedo, conversar, criança, cuidado, curiosidade, etário, gesto, infantil, infância, instituição, lúdico, mão, nascimento, ouvir, par, parceiro, percepção, sensorial, som
\end{textsample}
